\section*{Chapitre \ref{ch:iv2:behavioural-and-fit} --- Comportement et adéquation}
\markboth{\protect\omTocMark{}{Réponses~: Comportement et adéquation}}{\protect\omTocMark{}{Réponses~: Comportement et adéquation}}

\begin{solution}{iq:iv2:behavioural-and-fit:1}
Raconter une courte histoire qui mène à ce poste~: ce que l'on a étudié ou fait, une ou deux réalisations qui montrent qu'on sait faire ce travail, et pourquoi on est là aujourd'hui. S'arrêter vers quatre-vingt-dix secondes, et laisser le recruteur choisir ce qu'il veut approfondir.

Pour construire la sienne, rédiger trois phrases au passé, deux de preuves et une au présent. Choisir des preuves chiffrées, dont on est prêt à discuter pendant dix minutes.

\textbf{Illustration, avec les faits de chacun entre crochets.} «~J'ai étudié [matière] à [lieu], où j'ai découvert que j'aimais [type précis de problème]. J'ai ensuite travaillé sur [projet], où j'ai [ce que j'ai fait], et le résultat a été [chiffre]. Par ailleurs, j'ai construit [second élément], ce qui m'a appris [compétence]. Je postule à ce poste parce que [raison précise liée au travail]. Je suis prêt à approfondir chacun de ces points.~»

\iqwrong{Dérouler son CV dans l'ordre, ou commencer par son enfance. Cela prend cinq minutes, et n'apprend rien qu'il ne puisse déjà lire.}
\end{solution}

\begin{solution}{iq:iv2:behavioural-and-fit:2}
Nommer quelque chose de précis que ce travail a et que l'autre n'a pas~: par exemple, que le résultat de son travail se mesure vite et en argent, que son code ou son modèle touche de près à une décision, ou que le problème lui-même intéresse. Ne pas dire que cela tient à la rémunération, et ne pas prétendre que la rémunération ne compte pas.

Pour construire la sienne, lister ce qu'on a le plus apprécié dans son travail réel, et trouver l'élément que la finance y ajoute. Puis réfléchir à ce qu'on abandonne~: des produits que beaucoup de gens voient et utilisent, des travaux qu'on peut publier, et la liberté de partager ce qu'on apprend. Dire qu'on en est conscient.

\textbf{Illustration, avec les faits de chacun entre crochets.} «~Dans [projet], ce que j'ai le plus aimé, c'était [obtenir rapidement une réponse sur le fait d'avoir eu raison]. En finance, cette boucle est courte et difficile à contester, parce que [votre raison]. Je comprends qu'on ne peut souvent pas publier, et qu'une bonne idée cesse de fonctionner. Je l'accepterais pour [raison personnelle].~»

\iqwrong{«~J'ai toujours été passionné par les marchés~», sans rien derrière. On demandera ce que l'on a lu, négocié ou construit. Avoir une réponse prête.}
\iqfollow{Est-ce une question d'argent~?}{Être honnête~: la rémunération compte, et ce n'est pas la seule raison. Un démenti catégorique est moins crédible, et il lance mal l'échange.}
\end{solution}

\begin{solution}{iq:iv2:behavioural-and-fit:3}
Dire quelle partie du travail on veut faire toute la journée, et donner des preuves qu'on la pratique déjà par choix. Montrer qu'on respecte l'autre rôle, et ne pas le présenter comme un second choix.

En gros, la différence à examiner oppose la rapidité à la profondeur. Le trading récompense les décisions rapides dans l'incertitude, avec un résultat visible le jour même. La recherche récompense un travail long et soigneux, où la plupart des idées échouent et où le résultat arrive au bout de plusieurs mois. Ce que l'on préfère doit transparaître dans ce qu'on a déjà fait.

\textbf{Illustration pour un trader, avec les faits de chacun entre crochets.} «~Quand j'ai [jeu ou position gérée, avec ses propres mots], j'aimais décider avec une information incomplète et voir le résultat rapidement. Sur [projet long], j'ai constaté que j'étais moins patient à l'[étape finale]. Je préfère donc être au plus près de la décision, et travailler avec des chercheurs qui ont plus de patience pour la démonstration que moi.~»

\iqwrong{«~Je ne sais pas, je pourrais faire l'un comme l'autre~». On croirait qu'on ne s'est penché sur aucun des deux. Si l'on n'a vraiment aucune préférence, dire ce qui trancherait, et qu'on veut le savoir en faisant une chose précise ici.}
\end{solution}

\begin{solution}{iq:iv2:behavioural-and-fit:4}
Comparer la nature du travail plutôt que de classer les sociétés~: comment on gagne de l'argent, à quel point on en est proche, à quelle vitesse on sait si l'on a eu raison, et quels sont les risques. Donner ensuite la raison qui correspond à son profil.

Pour construire la sienne, préparer une phrase par type, en s'appuyant sur ses lectures, et vérifier chacune auprès de sources publiques. En gros~: un teneur de marché gagne de petites sommes répétées en cotant des prix et en gérant le risque, un hedge fund garde ses positions plus longtemps, une banque sert des clients et gère des risques réglementés, et une société crypto évolue sur des marchés plus jeunes, avec davantage de risque opérationnel. Le prendre comme point de départ, car les sociétés diffèrent beaucoup.

\textbf{Illustration, avec les faits de chacun entre crochets.} «~Je veux [type] parce que je veux [caractéristique], et j'en ai eu la preuve quand j'ai [preuve]. J'ai regardé [autre type], et je regretterais [caractéristique]. Il s'agit d'une préférence, pas d'un jugement sur eux.~»

\iqwrong{Critiquer les trois autres. On peut très bien parler à quelqu'un qui a travaillé dans l'une d'elles, et cela en dit plus sur soi que sur elles.}
\end{solution}

\begin{solution}{iq:iv2:behavioural-and-fit:5}
Donner un bref exposé précis, fondé sur des sources publiques~: les instruments que la société négocie, sur quels marchés, comment elle gagne de l'argent, et ce qui pourrait lui nuire. Dire ensuite ce que l'on ignore, et poser la question.

Pour se préparer, lire les pages de la société, ses travaux et interventions publiés, les dépôts publics éventuels et l'actualité récente. Noter trois faits, une question et un risque. Le risque doit être réel~: pour une société qui affiche des cotations, par exemple, une baisse des volumes, des concurrents plus rapides ou une panne technique, chacun pouvant se démontrer à partir de la façon dont une telle société gagne sa vie.

\textbf{Illustration, avec les faits de chacun entre crochets.} «~D'après [source publique], je comprends que la société fait surtout [activité] sur [marchés], donc elle gagne quand [mécanisme]. Je pense que ses principaux risques sont [risque un] et [risque deux]. Je n'ai pas pu trouver [élément précis]. Est-ce un sujet que l'on peut aborder~?~»

\iqwrong{Répéter la page d'accueil, ou présenter une rumeur lue sur un forum comme un fait. Les rumeurs sur les rémunérations et les méthodes sont souvent périmées, et parfois fausses. Si l'on en cite une, dire d'où elle vient et qu'on ne peut pas s'en porter garant.}
\end{solution}

\begin{solution}{iq:iv2:behavioural-and-fit:6}
Nommer une vraie faiblesse qui compte un peu pour ce poste, donner un exemple, et dire ce qu'on fait pour y remédier et dans quelle mesure cela a fonctionné.

Pour en trouver une, se demander ce qu'un collègue ou un enseignant a récemment demandé d'améliorer, ou ce qui a mal tourné dans le projet pendant sa pire semaine. Ne pas choisir un défaut fatal pour le poste. Un trader qui affirme se figer sous la pression n'a pas répondu à la bonne question.

\textbf{Illustration, avec les faits de chacun entre crochets.} «~J'ai tendance à [habitude réelle, par exemple peaufiner trop longtemps avant de montrer le travail]. Sur [projet], cela a coûté [effet concret]. Maintenant, je [changement précis, par exemple montrer une version brute dès le deuxième jour], et cela a aidé, même si je [ce qui reste]~».

\iqwrong{«~Je suis perfectionniste~», ou «~Je travaille trop~». On y entend tout de suite une esquive. Cela laisse aussi entendre qu'on ne donnera pas de réponse franche face à un point gênant, ce qui, sur un desk, est la faiblesse qui compte le plus.}
\end{solution}

\begin{solution}{iq:iv2:behavioural-and-fit:7}
Décrire le travail que l'on veut faire et le niveau de maîtrise visé. Ne pas donner de titre, ni de projet de départ. Ajouter ce que l'on aurait appris d'ici là, et dire qu'on peut se tromper sur les détails.

Pour construire le sien, se renseigner sur ce que fait réellement quelqu'un qui est en poste depuis cinq ans dans ce type de société, et lui poser la question en entretien si possible. Le plan doit montrer qu'on veut ici de la profondeur. Son échelle de temps doit correspondre à celle du poste.

\textbf{Illustration, avec les faits de chacun entre crochets.} «~J'aimerais être quelqu'un qui est responsable de [type de travail précis] et qui sait [compétence qui me manque aujourd'hui]. Pour y arriver, je pense consacrer les premières années à apprendre [quoi]. Je ne sais pas encore si cela passe par [voie un] ou par [voie deux], et c'est l'une des choses que je veux découvrir ici~».

\iqwrong{«~Diriger mon propre fonds~», ou «~à la place du recruteur~». L'ambition est une bonne chose, mais cela laisse entendre qu'on voit le poste comme un tremplin. Une réponse du type «~Je n'y ai pas réfléchi~» est tout aussi faible.}
\end{solution}

\begin{solution}{iq:iv2:behavioural-and-fit:8}
Ne pas se comparer à des personnes qu'on n'a jamais rencontrées. Donner les deux ou trois apports que l'on fait, chacun appuyé par une preuve, et les relier à ce qu'exige ce poste.

Pour construire la sienne, lire la description du poste et en choisir la partie la plus difficile. Puis choisir quelque chose que l'on a fait et que la plupart des gens au même parcours n'ont pas fait~: par exemple, une réalisation que l'on a bâtie, un résultat que l'on peut montrer, ou une combinaison de compétences peu courante. Le dire simplement, puis s'arrêter. Une bonne réponse prend moins d'une minute.

\textbf{Illustration, avec les faits de chacun entre crochets.} «~Je ne peux pas parler des autres. Ce que je peux dire, c'est que j'ai [preuve un, avec un chiffre] et [preuve deux]. Pour ce poste, cela veut dire que je pourrais [apport concret] plus vite que quelqu'un qui devrait l'apprendre. Et je travaille bien avec [type de personnes], par exemple lorsque [exemple court et vrai]~».

\iqwrong{Être si modeste qu'on n'apporte rien, ou se déclarer le meilleur. On a demandé de défendre un dossier, et ni l'un ni l'autre n'en est un.}
\end{solution}

\begin{solution}{iq:iv2:behavioural-and-fit:9}
Voici une version d'une minute que l'on peut adapter. «~Prenons un magasin qui achète et revend des voitures d'occasion. Il rachète une voiture 9\,000 et en revend une semblable 9\,500, et il gagne les 500 en faisant cela toute la journée, avec quiconque franchit la porte. Il n'a pas besoin de savoir si les prix des voitures vont monter. Un teneur de marché fait de même avec des actions et d'autres instruments, mais sur écran, par petites quantités, plusieurs fois par seconde. Il perd de l'argent de deux façons. Si les prix baissent alors qu'il détient un gros inventaire, il supporte la perte. Et si le vendeur sait quelque chose qu'il ignore, comme un défaut de la voiture, il a payé trop cher. Il garde donc un inventaire réduit et ajuste ses prix rapidement.~»

Pour construire la sienne, choisir une image que l'auditeur connaît déjà, et la tester auprès de quelqu'un qui n'est pas dans la finance.

\iqwrong{Employer les mots fourchette, liquidité, inventaire ou sélection adverse. Si l'on a besoin de ces mots, on n'a pas expliqué. Ne pas oublier la partie sur les pertes, qui est en réalité l'objet de la question.}
\iqfollow{Alors pourquoi les gens négocient-ils avec lui, s'il prélève toujours sa marge~?}{Parce qu'ils peuvent négocier tout de suite, à tout moment, pour un coût réduit. Sans lui, ils pourraient attendre longtemps avant de trouver un acheteur.}
\end{solution}

\begin{solution}{iq:iv2:behavioural-and-fit:10}
Dire ce qui s'est passé en une ou deux phrases simples, assumer sa part de responsabilité, dire ce qui a changé, et donner une preuve que cela a bien changé. Puis s'arrêter. Ne pas s'excuser à répétition, et ne pas raconter une longue histoire.

Pour construire la sienne, décider de ce qui est vrai et de ce que l'on veut partager. Il n'est pas nécessaire de donner des détails privés, comme un problème de santé ou de famille~: «~une question de santé familiale, aujourd'hui réglée~» suffit. Les preuves doivent venir après le point faible~: un résultat ultérieur, un projet achevé.

\textbf{Illustration, avec les faits de chacun entre crochets.} «~En [année], mes résultats sont tombés à [résultat]. La raison était [raison courte et vraie]. J'ai changé [ce que j'ai fait], et ensuite j'ai [résultat ultérieur]. La question est légitime, et je pense que [résultat ultérieur] est un meilleur indicateur de ce que je peux faire.~»

\iqwrong{Rejeter la faute sur un enseignant, un responsable ou un cours, ou modifier les dates de son CV. Une vérification peut le révéler, et cela mettrait fin au processus. Cacher une mauvaise nouvelle est ce qui inquiète le plus un employeur.}
\iqfollow{Comment être sûr que cela ne se reproduira pas~?}{Dire ce qui a changé depuis, et donner un moment ultérieur où l'on a subi la même pression.}
\end{solution}

\begin{solution}{iq:iv2:behavioural-and-fit:11}
Être honnête sur le fait qu'on passe des entretiens ailleurs, dire globalement où, et dire ce qui trancherait. Il n'est pas obligatoire de donner des noms, et il ne faut jamais inventer une offre.

Si l'on a une échéance, le dire tôt et calmement~: «~J'ai une offre qui expire le [date], il me serait donc utile de savoir où j'en suis d'ici là~». C'est en général préférable pour les deux camps à une surprise de dernière minute.

\textbf{Illustration, avec les faits de chacun entre crochets.} «~Oui, je discute avec [nombre] autres sociétés de [type]. Ce qui tranchera pour moi, c'est [ce qui compte pour soi]. D'après ce que j'ai vu, cette société est forte sur [point vrai], et je suis encore en train de me renseigner sur [autre point]~».

\iqwrong{«~Votre société est mon seul choix.~» Elle est rarement crue, et si elle est fausse, la première déclaration faite devant eux est un mensonge. L'autre erreur consiste à faire jouer les sociétés les unes contre les autres avec une offre que l'on n'a pas. C'est un petit monde, et les gens en parlent.}
\end{solution}

\begin{solution}{iq:iv2:behavioural-and-fit:12}
Dire ce que l'on fait réellement, et ce que cela a appris ou n'a pas appris. Les jeux peuvent enseigner à raisonner en cotes, à dimensionner ses mises et à perdre avec calme. Ils ne montrent pas que l'on sait trader. Si l'on ne joue pas, il faut le dire et donner ce que l'on possède.

Pour élaborer la sienne, il faut être prêt à donner un chiffre~: comment on sait si l'on gagne. Une personne qui joue de l'argent sans pouvoir dire son résultat à long terme, ni comment elle l'a consigné, n'a pas tiré la principale leçon.

\textbf{Illustration, avec ses propres faits entre crochets.} «~Je joue à [jeu] [fréquence], et j'ai tenu un registre sur [période], qui montrait [résultat]. Cela m'a appris [leçon sur le dimensionnement ou le tilt, dans ses propres mots]. En quoi cela diffère du trading~: [différence réelle].~» Si l'on ne joue pas~: «~Je ne joue pas. Là où j'ai le plus réfléchi aux cotes et aux mauvais résultats, c'est [votre exemple].~»

\iqwrong{Raconter l'histoire d'un gros gain, ou laisser entendre qu'une habitude de jeu revient à faire du trading. Quelqu'un qui parie sans avantage, ou pour ressentir quelque chose, est exactement ce que redoute un responsable du risque.}
\end{solution}

\begin{solution}{iq:iv2:behavioural-and-fit:13}
Choisir un sujet difficile et qui ne fait pas partie de ce que l'on a étudié. Dire comment on a choisi ses sources, comment on s'est entraîné, et quel test a montré que l'on avait appris. La dernière partie est celle que la plupart des gens omettent.

Pour élaborer la sienne, il faut décrire la difficulté avec honnêteté~: ce que l'on a d'abord mal compris, et comment on s'en est aperçu. «~J'ai lu un livre~» ne constitue pas un récit d'apprentissage. Une preuve peut être un programme qui fonctionne, un ensemble d'exercices résolus, une explication acceptée par un expert, ou un résultat vérifié contre une réponse connue.

\textbf{Illustration, avec ses propres faits entre crochets.} «~Je me suis formé seul à [sujet] pendant [durée], parce que [raison]. J'ai commencé par [source], et je me suis bloqué sur [point]. Alors j'ai [ce que l'on a fait]. J'ai su que cela avait marché quand [test, par exemple~: mon code reproduisait un chiffre publié, ou j'ai résolu un ensemble d'exercices sans aide].~»

\iqwrong{Choisir quelque chose de facile, ou tiré d'un cours, parce que cela paraît sûr. Choisir un sujet dont on ne peut pas discuter en profondeur est pire, car les recruteurs peuvent le connaître mieux que le candidat.}
\end{solution}

\begin{solution}{iq:iv2:behavioural-and-fit:14}
Poser deux ou trois vraies questions auxquelles on n'aurait pas pu répondre à partir de sources publiques, portant sur le travail de la personne et sur la manière dont la réussite est jugée. Réserver les questions sur la rémunération et les conditions à la période qui suit une offre.

Pour un futur manager, demander quelles erreurs commet le plus souvent une nouvelle recrue pendant les six premiers mois, comment on serait jugé au bout d'un an, sur quel sujet difficile travaille l'équipe, et comment les désaccords sont tranchés. Pour un associé senior, demander ce qui a changé dans l'activité au cours des dernières années, ce qui l'inquiète, et ce qu'il recherche chez les personnes qui restent. Après une offre, demander la rémunération et son mode de fixation, les préavis et toute restriction sur le fait de travailler ailleurs après le départ, quelle part du code ou de la recherche on peut réutiliser ailleurs, et comment une position ou un projet est transféré à quelqu'un d'autre en cas d'échec.

\iqwrong{«~Non, on a tout couvert.~» Cela laisse penser à un faible intérêt. Poser une question dont la réponse figure sur le site, ou demander d'abord les congés ou le montant du bonus, est aussi un mauvais signal.}
\iqfollow{Et s'ils vous répondent de façon vague~?}{Poser une relance, poliment, et la noter. Une réponse vague sur la manière dont on juge les gens est une information dont on a besoin pour sa décision.}
\end{solution}

\begin{solution}{iq:iv2:behavioural-and-fit:15}
Choisir une vraie déception dont on a assumé la responsabilité et que l'on peut expliquer avec des chiffres. Dire ce que l'on attendait et pourquoi, ce qui s'est passé, pourquoi l'on pense que c'est arrivé et comment on l'a vérifié, et ce que l'on fait désormais différemment, concrètement.

Pour élaborer la sienne, il faut utiliser quatre parties~: attente, résultat, cause, changement. Consacrer l'essentiel du temps à la cause et au changement. Les meilleures réponses distinguent un mauvais processus de la malchance. Parfois, on a pris la bonne décision avec ce que l'on savait. Il faut le dire, et expliquer comment on le saurait.

\textbf{Illustration, avec ses propres faits entre crochets.} «~Sur [projet], j'attendais [résultat chiffré], parce que [raison]. Le résultat s'est établi à [nombre]. En cherchant à comprendre, la cause était [cause], que j'ai vérifiée par [vérification]. Une part en tenait à la malchance, parce que [raison], mais la part dont je suis responsable est [son erreur]. Alors je fais désormais [changement précis], que j'ai aussi appliqué sur [exemple ultérieur].~»

\iqwrong{Choisir un succès déguisé («~nous n'avons dépassé l'objectif que de 5~\%~»), ou rejeter la faute sur l'équipe, les données ou le marché. Aucune de ces solutions ne laisse quoi que ce soit à en tirer.}
\iqfollow{Qu'avez-vous dit aux autres lorsque vous avez vu que cela tournait mal~?}{Avoir en tête une date~: celle où l'on a su, et celle où on l'a dit. Un long écart entre ces deux dates est ce que les recruteurs recherchent.}
\end{solution}

\begin{solution}{iq:iv2:behavioural-and-fit:16}
Choisir une vraie erreur, au coût réel, que l'on a repérée ou reconnue soi-même, et se concentrer sur ce qui s'est passé après sa découverte~: à qui on l'a signalée, en combien de temps, ce que l'on a fait pour la corriger et le contrôle ajouté pour qu'elle ne se reproduise pas.

Pour élaborer la sienne, il faut choisir un cas à la mesure de son niveau, que l'on puisse raconter en deux minutes environ. Éviter tout ce qui nuit à un ancien collègue ou qui trahit une confidence.

\textbf{Illustration, avec ses propres faits entre crochets.} «~Dans [contexte], j'ai commis [erreur], et le coût a été de [chiffre ou effet]. Je l'ai repérée quand [comment]. Dans un délai de [durée], j'ai prévenu [personne], et nous avons [correction]. Je fais désormais [contrôle ou habitude], pour que le même problème soit détecté par [comment].~»

\iqwrong{Une erreur trop minime pour compter («~j'ai une fois oublié une virgule~»), ou une erreur sans action ensuite. Pire encore est un long intervalle entre la découverte et le moment où l'on en parle à quelqu'un. L'erreur est pardonnable, mais le retard l'est souvent moins.}
\iqfollow{Que se serait-il passé si vous n'aviez rien dit~?}{Répondre avec franchise, même si la réponse est «~peut-être que personne ne l'aurait jamais su~». On peut alors expliquer pourquoi on l'a dit malgré tout.}
\iqwhy{Les desks peuvent survivre aux erreurs signalées immédiatement. Ce qui fait du mal aux sociétés, c'est qu'une personne cache une petite perte jusqu'à ce qu'elle devienne importante. C'est la réponse dont le recruteur se souviendra.}
\end{solution}

\begin{solution}{iq:iv2:behavioural-and-fit:17}
Établir d'abord les faits, puis prévenir la personne qui a décidé, corriger ensuite le rapport, puis traiter la cause. Le faire aujourd'hui même. Il faut le dire, même si personne ne l'aurait jamais remarqué.

Dans l'ordre~: d'abord, vérifier rapidement que l'erreur est bien réelle et mesurer son ampleur. Y consacrer quelques minutes ou une heure au plus. Ensuite, prévenir directement le responsable et la personne qui a utilisé le chiffre, le jour même~: ce qui était faux, le chiffre exact et ce qu'il change dans leur décision. Troisièmement, republier le rapport avec la correction signalée. Quatrièmement, chercher pourquoi l'erreur s'est produite, et ajouter un test ou un contrôle. Cinquièmement, vérifier si la même erreur figure dans les autres documents envoyés.

\textbf{Illustration des formulations.} «~Le chiffre de [chiffre faux] à la page [page] du rapport de la semaine dernière est faux. Le bon chiffre est [bon chiffre]. Il vient de mon erreur à l'étape [étape]. Pour la décision sur [décision], cela signifie [effet]. Je corrige le rapport maintenant, et je vérifie [quoi d'autre].~»

\iqwrong{Attendre d'avoir une analyse complète, ou un bon moment, ou corriger discrètement la fois suivante. Cela transforme une petite erreur en une question sur soi.}
\iqfollow{Et si votre responsable vous dit de la laisser, puisque cela ne change pas grand-chose~?}{Dire pourquoi on préfère la corriger, et le consigner par écrit. Si cela touche un client ou un chiffre réglementaire, il faut la faire remonter par la voie appropriée.}
\end{solution}

\begin{solution}{iq:iv2:behavioural-and-fit:18}
Raconter une situation dans laquelle on a soulevé le problème avec des preuves, en privé, en écoutant, et où le résultat final a été tranché par les faits ou par une décision claire. Dire simplement qui s'est révélé avoir raison, parce que la question est posée. La manière dont on s'est comporté compte davantage que le score, tout comme le fait de s'incliner devant la décision.

Pour élaborer la sienne, il faut utiliser cinq parties~: ce que chacun pensait, ce que l'on a apporté, comment l'affaire a été tranchée, qui s'est révélé avoir raison, et ce que l'on a appris sur la manière de travailler avec cette personne. Un récit dans lequel on a eu tort et où on le reconnaît peut être plus convaincant qu'un récit où l'on a gagné.

\textbf{Illustration des formulations.} «~[Personne senior] voulait [approche]. Je pensais que [autre approche], parce que [preuve]. J'ai demandé [délai], et je leur ai montré [test ou données]. Ils [résultat]. Ensuite [ce qui s'est passé, et qui avait raison]. J'ai appris [leçon vraie].~»

\iqwrong{Un récit qui se termine par «~et j'avais raison~», avec un imbécile en face. L'autre réponse faible est «~je ne suis jamais en désaccord avec personne~», qui laisse entendre qu'on ne prendra pas la parole lorsqu'une position est erronée.}
\end{solution}

\begin{solution}{iq:iv2:behavioural-and-fit:19}
Choisir un cas où quelque chose semblait trop beau ou simplement étrange, et reprendre ce qui a fait naître le doute, comment on l'a testé et quelle en était la réponse. Découvrir un bogue dans son propre bon résultat est la version la plus forte.

Pour élaborer la sienne, il faut utiliser~: le signal qui indiquait que quelque chose clochait, la liste des causes possibles, le test le moins coûteux pour chacune, et la conclusion. Parmi les bons contrôles figurent un résultat connu sur de petites données, un changement de période, le retrait des principaux contributeurs, la permutation des étiquettes et une réimplémentation indépendante.

\textbf{Illustration des formulations.} «~[Modèle ou analyse] a donné [résultat], ce qui semblait [trop beau ou étrange] parce que [référence]. Je soupçonnais [cause un] ou [cause deux]. J'ai testé par [contrôle], ce qui a montré [constat]. Le chiffre exact était [nombre]. J'ai prévenu [personne], et nous avons [action].~»

\iqwrong{Choisir un cas où l'on a douté d'un résultat seulement parce qu'il déplaisait. Le doute doit reposer sur une raison que l'on peut énoncer, et être tout aussi prêt à douter d'un résultat qui plaît. L'autre erreur consiste simplement à relancer le modèle et à l'accepter.}
\end{solution}

\begin{solution}{iq:iv2:behavioural-and-fit:20}
Leur parler d'abord, en privé, rapidement et avec des faits précis. Proposer de l'aide, et convenir de ce qu'ils livreront et de la date. Si cela ne suffit pas, ajuster le plan avec l'équipe. Prévenir un supérieur dès que l'échéance est réellement menacée, sans surprise.

L'ordre est donc le suivant~: d'abord, chercher la cause. La personne peut être bloquée, malade, ou ne pas savoir clairement ce qui est attendu. Ensuite, découper le travail en petites tâches visibles, avec une date pour chaque partie. Troisièmement, si ces échéances ne sont pas tenues, lui dire clairement qu'il faut en parler, puis en parler. Énoncer des faits sur le travail, et non sur la personne. Entre-temps, protéger l'échéance en faisant les parties les plus importantes, et dire ce qui ne sera pas fait.

\textbf{Illustration des formulations.} «~J'ai remarqué que [partie] n'est pas terminée, et nous en avons besoin pour [date]. Y a-t-il un obstacle~? Est-il possible d'avoir [petite partie] pour [date]~? Si cela pose problème, il faut le dire maintenant, et l'on peut changer le plan.~»

\iqwrong{Faire en silence sa part en en gardant de la rancœur, ou aller tout de suite voir le supérieur. Le premier cache un problème à la personne qui doit le savoir. Le second met fin à la relation de travail avant même d'avoir posé la question.}
\end{solution}

\begin{solution}{iq:iv2:behavioural-and-fit:21}
Ne pas le mettre en production sans test de sa propre initiative, et ne pas se contenter de refuser. Savoir ce qui est en jeu, proposer la voie la plus rapide et sûre, et faire décider à la personne investie de l'autorité si le risque doit être pris.

Les étapes~: d'abord, demander quel est le changement, ce qu'il rapporterait et ce qui se passe s'il n'est pas en production aujourd'hui. Ensuite, dimensionner le risque~: quel est le pire scénario, et quelles limites ou quels interrupteurs permettront de l'arrêter~? Troisièmement, proposer des options, comme une version plus réduite, un passage en shadow mode, uniquement sur une petite position ou un seul symbole, derrière un feature flag, ou avec une limite de taille, et des tests limités à ce que permet l'heure. Quatrièmement, si le processus de la société prévoit une façon d'approuver un changement d'urgence, s'en servir, et obtenir l'approbation par écrit, de la part de celui qui porte le risque.

\textbf{Illustration des formulations.} «~Je ne peux pas le mettre en production à l'aveugle, car si cela tourne mal la perte serait de [pire cas]. Je peux avoir prête [version réduite et sûre] dans [délai], avec une limite de [taille]. Si la version complète est souhaitée aujourd'hui, il faut que [responsable nommé] l'approuve, et je la préparerai pour pouvoir la retirer en [délai].~»

\iqwrong{«~Bien sûr~», parce que le trader est senior et pressé. Tout aussi faux est un refus sec, sans alternative. Une mise en production précipitée peut faire perdre très vite beaucoup d'argent.}
\iqfollow{Et si cela tourne mal~?}{On la coupe d'abord, on prévient ensuite tout le monde, et seulement après on cherche la cause. Prévoir le moyen de la couper avant de la mettre en production.}
\iqwhy{Ici, un ingénieur est jugé à la fois sur son aide au desk et sur le fait de ne pas faire tomber la société. C'est le moment où ces deux exigences se rencontrent.}
\end{solution}

\begin{solution}{iq:iv2:behavioural-and-fit:22}
Choisir un vrai échec dont on a assumé la responsabilité. Le raconter dans l'ordre chronologique~: comment on l'a repéré, ce qui a été fait pour arrêter les dégâts, à qui on l'a signalé, comment on a trouvé la cause, et ce qui a été changé pour que cela soit détecté plus tôt. Arrêter d'abord les dégâts, et chercher la cause ensuite.

Pour en construire un, il faut rédiger un bref post-mortem\index{post-mortem} comme s'il était destiné à l'équipe~: impact, chronologie, cause, ce qui a bien fonctionné et actions menées. Donner un chiffre pour l'impact, ainsi que pour le délai de détection et le délai de rétablissement. Être précis sur sa propre part.

\textbf{Illustration, avec ses propres faits entre crochets.} «~Le [date], [système] a subi [défaillance]. Nous l'avons découverte grâce à [alerte ou personne] au bout de [durée]. Dans les dix premières minutes, j'ai [action pour arrêter les dégâts] et prévenu [personne]. La cause s'est révélée être [cause]. L'impact a été de [chiffre]. Ensuite, nous avons ajouté [test, alerte ou limite], de sorte que cela soit désormais détecté en [délai].~»

\iqwrong{Raconter une nuit blanche héroïque, alors que l'essentiel est que la défaillance n'aurait jamais dû exiger une telle nuit. En rejeter la faute sur quelqu'un d'autre est pire. Si l'on n'a jamais connu de défaillance, c'est probablement que l'on n'a pas beaucoup assumé de responsabilités.}
\end{solution}

\begin{solution}{iq:iv2:behavioural-and-fit:23}
Choisir une vraie échéance où il a fallu abandonner quelque chose. Montrer comment on a décidé~: ce que l'on a classé comme essentiel, à qui l'on a annoncé ce qui ne serait pas fait, et comment le résultat s'est dessiné. La décision est la partie intéressante, et non les longues heures.

Pour en construire un, retenir~: quelle était l'échéance et pourquoi elle ne pouvait pas être repoussée, les options, la manière dont on a choisi, à qui on l'a annoncé, et le résultat. Inclure ce que l'on avait mal évalué dans la première estimation. Laisser ses heures de côté, car les longues heures ne prouvent rien.

\textbf{Illustration, avec ses propres faits entre crochets.} «~Nous disposions de [durée] pour livrer [élément] parce que [raison]. À [moment], j'ai vu que nous ne pouvions pas tout faire. J'ai classé [parties] selon [critère], abandonné [partie], et prévenu [personne] dès le matin. Nous avons livré [résultat], et terminé [partie] avant le [date]. La prochaine fois, je [ce que je ferais autrement].~»

\iqwrong{«~J'ai simplement travaillé toute la nuit, et tout fait.~» Cela élude la question, et laisse penser qu'on ne signale peut-être pas ses difficultés.}
\end{solution}

\begin{solution}{iq:iv2:behavioural-and-fit:24}
Pour la plupart des gens, prendre le million, même si le pile ou face rapporte en moyenne 1.5 million. Avec une utilité logarithmique, on est indifférent à une richesse initiale d'exactement 1 million. En dessous, on prend la somme certaine, et au-dessus, on prend le pile ou face.

Avec une richesse $W$ en millions, l'indifférence s'écrit $\ln(W + 1) = \tfrac12\ln(W + 3) + \tfrac12\ln W$, donc $(W + 1)^2 = W(W + 3)$, ce qui donne $W = 1$. Ainsi, quelqu'un qui possède 100\,000 dollars doit prendre la somme certaine, et quelqu'un qui possède 10 millions doit prendre le pile ou face. La réponse d'une personne doit dépendre de sa richesse, et il faut le dire.

Ainsi, la prime abandonnée, 0.5 million en espérance, sert à ne pas risquer une grande partie de sa richesse. Pour $W = 10$, l'équivalent certain\index{equivalent certain@équivalent certain} du pile ou face vaut environ 1.40 million, plus proche de sa moyenne. Pour $W = 0.1$, il vaut environ 0.46 million.

\iqwrong{«~Je prendrais le pile ou face, car l'espérance est plus élevée~», sans réfléchir à sa richesse. Un recruteur peut enchaîner en vous demandant de miser votre propre maison. L'erreur inverse consiste à traiter de la même façon un petit pari, par exemple à refuser un petit pari favorable simplement parce que c'est un pari.}
\iqfollow{Que doit faire une société disposant de 1 milliard~?}{Prendre le pile ou face, et tous les paris semblables, car chacun est petit par rapport à son capital. C'est pourquoi une société peut prendre des risques qu'une personne ne pourrait pas prendre.}
\iqwhy{Cela montre si l'on distingue l'espérance de ce que l'on peut se permettre de perdre, idée à la base de tout dimensionnement.}
\end{solution}

\begin{solution}{iq:iv2:behavioural-and-fit:25}
Choisir un vrai risque, et non une fanfaronnade financière. Dire quelles étaient à l'époque ses estimations de chances et de coûts, ce que l'on a fait pour limiter le revers possible, et juger la décision indépendamment du résultat. Dire si on la referait avec ce que l'on savait alors.

Pour en construire un, peuvent convenir un déménagement, un changement de domaine, le départ d'un emploi sûr, la création d'une activité ou un pari sur une compétence. Il faut que l'échec ait coûté quelque chose de réel. Énoncer son estimation d'alors, même si elle était approximative. Donner ensuite la partie la plus intéressante~: ce que l'on a appris sur la qualité de cette estimation.

\textbf{Illustration, avec ses propres faits entre crochets.} «~J'ai [décision] en [année]. Je pensais avoir environ [chance] de [bon résultat], et le revers était [coût], que je pouvais supporter parce que [raison]. Il s'est avéré [résultat]. Je [le ferais ou ne le ferais pas] à nouveau avec ce que je savais, parce que [raison]. Ce que j'ai mal évalué, c'est [quoi].~»

\iqwrong{Dire qu'une décision était bonne parce qu'elle a fonctionné, ou présenter une démarche imprudente comme du courage. Quelqu'un qui prend un risque important sans aucune estimation des chances est précisément ce qu'une société ne veut pas.}
\end{solution}

\begin{solution}{iq:iv2:behavioural-and-fit:26}
Les deux stratégies rapportent en moyenne 0.1 par jour, et $B$ est la meilleure~: son écart type vaut environ la moitié de celui de $A$, si bien que son rapport entre moyenne et écart type est environ deux fois meilleur. Laquelle des deux est la plus difficile à supporter dépend de chacun. $A$ subit de rares pertes importantes, et $B$ enchaîne de longues séries de pertes ternes.

\begin{center}\small
\begin{tabular}{@{}lrr@{}}
 & $A$ & $B$ \\ \hline
Moyenne par jour & 0.10 & 0.10 \\
Écart type par jour & 2.70 & 1.35 \\
Moyenne sur 250 jours & 25 & 25 \\
Écart type sur 250 jours & 42.7 & 21.3 \\
Moyenne rapportée à l'écart type & 0.037 & 0.074 \\
Moyenne annuelle rapportée à l'écart type, 250 jours & 0.59 & 1.17 \\
Plus longue série typique de jours perdants sur 250 & 2 & 9 \\
Pire perte sur un jour & 8 & 1
\end{tabular}
\end{center}
Avec $A$, on a une bonne journée 90~\% du temps, puis une perte de 8 efface huit bonnes journées. Avec $B$, on perd la plupart des jours, et une année typique comporte une série d'environ 9 pertes consécutives, ce qui explique que beaucoup de gens l'abandonnent au pire moment. Les longueurs de séries proviennent d'un décompte exact de la plus longue série médiane sur 250 jours.

Une réponse solide dit que $B$ est meilleure selon les chiffres, que $A$ paraît meilleure tant qu'on ne voit pas la queue de distribution, et comment savoir laquelle on pourrait continuer à suivre. Avant de répondre, se demander si la queue de $A$ ne vaut vraiment que 8.

\iqwrong{«~$A$, car elle gagne 90~\% du temps.~» La fréquence de gain ne dit rien en soi, et un taux de réussite (hit rate) élevé est souvent l'allure que prend un vendeur de risque de queue.}
\iqfollow{Que vérifieriez-vous avant de faire confiance aux chiffres de $A$~?}{La longueur de l'historique. Un taux de pertes de 10~\% se voit vite, mais une perte bien plus rare et plus importante n'a peut-être pas encore eu lieu.}
\iqwhy{Le tempérament doit correspondre au profil de résultats de la stratégie. Les gens abandonnent de bonnes stratégies pendant les mauvaises séries normales.}
\end{solution}

\begin{solution}{iq:iv2:behavioural-and-fit:27}
Arrêter d'ajouter du risque, réduire la position comme l'exige la limite, prévenir immédiatement son responsable en lui donnant les faits, puis passer la journée à chercher si la perte vient de ses décisions ou d'un dysfonctionnement. Ne pas chercher à la rattraper.

L'ordre~: d'abord, respecter la règle. Une limite de perte existe précisément pour ce moment, donc clôturer ou réduire la position comme elle le prescrit. Ensuite, vérifier que la perte est réelle~: les cours de valorisation, les positions et les données sont-ils corrects~? Troisièmement, prévenir son responsable en chiffres~: combien, à cause de quoi, ce que l'on détient maintenant et ce que l'on prévoit. Quatrièmement, décomposer la perte en parties, par exemple le mouvement du marché, la taille de la position et l'exécution. Cinquièmement, consacrer le reste de la journée à un travail que l'on sait bien faire, comme l'analyse de la transaction, ou l'observation sans trader.

\textbf{Illustration des formulations.} «~Je suis en perte de [montant], ce qui atteint ma limite. Elle vient surtout de [position], parce que [cause]. Je l'ai [ramenée à telle taille]. Les cours de valorisation sont vérifiés. Je pense que [jugement]. Je n'ajouterai rien aujourd'hui, sauf si vous le souhaitez.~»

Pour le jour suivant, reprendre à une taille réduite, et ne la relever que lorsque le processus fonctionne à nouveau. Se demander si l'idée de départ tient toujours, en la traitant comme une nouvelle transaction.

\iqwrong{Trader plus fort pour rattraper la perte. C'est la façon classique de transformer une mauvaise journée en une journée bien pire. Cacher la taille de la position à son responsable est pire encore.}
\iqfollow{Et si vous êtes sûr que la position va remonter~?}{Tout le monde l'est toujours. Si l'on avait raison sur la position, le desk peut la reconstituer l'esprit clair. Une limite que l'on outrepasse dès que l'on est confiant n'est pas une limite.}
\iqwhy{Les mauvaises journées arrivent à tous les traders. Ce que la société ne peut pas accepter, c'est une personne dont le comportement se dégrade après une perte.}
\end{solution}

\begin{solution}{iq:iv2:behavioural-and-fit:28}
Traiter cela comme une infraction, et non comme un succès. Le dire soi-même à son responsable, établir ce que la règle aurait rapporté et quel était le risque supplémentaire, puis décider ouvertement si c'est la règle ou son comportement qui doit changer. Une règle enfreinte qui a rapporté est plus dangereuse qu'une règle enfreinte qui a fait perdre.

Le profit ne prouve pas que la décision était bonne. Séparer les éléments~: combien aurait rapporté la transaction à la taille prévue par la règle, et ce que la taille supplémentaire a ajouté~? Combien aurait pu perdre la taille supplémentaire si l'événement avait tourné dans l'autre sens~? Si cette perte aurait dépassé ce que l'on pouvait accepter, alors on a pris un risque que l'on avait décidé de ne pas prendre.

Ensuite, se demander pourquoi. Soit la règle est mauvaise, auquel cas on la change délibérément, par écrit, avec des preuves, à tête reposée. Soit la règle est juste, et on l'a enfreinte parce que l'on se sentait sûr de soi, et c'est ce sentiment qu'il faut surveiller. La prochaine fois, écrire la raison de dépasser la taille avant de trader, et demander à quelqu'un d'autre de la relire.

\textbf{Illustration des formulations.} «~Ma meilleure transaction de l'année a enfreint ma règle de dimensionnement. Voici ce que la règle aurait donné, et voici ce que j'aurais perdu si cela avait tourné dans l'autre sens. Je veux discuter de la question de savoir si la règle doit changer. Elle ne doit pas changer simplement parce que celle-ci a fonctionné.~»

\iqwrong{Garder discrètement le profit, ou se féliciter de son instinct. Un responsable qui l'apprend plus tard se demandera ce qui d'autre n'a pas été signalé.}
\iqfollow{Et si votre responsable est ravi et vous demande d'en faire davantage~?}{Lui montrer les mêmes chiffres, et demander une modification explicite de ses limites si c'est ce qu'il souhaite. Il ne faut pas en rester à une impression.}
\end{solution}

\begin{solution}{iq:iv2:behavioural-and-fit:29}
Vérifier discrètement et rapidement les faits, puis remonter le problème par la voie appropriée~: d'abord à son responsable, sauf s'il est impliqué, puis à la gestion des risques, au product control ou à la conformité. Ne pas l'ignorer. Ne pas accuser qui que ce soit avant d'avoir des faits, et ne pas confronter seul quelqu'un avec une accusation. Si rien ne se passe ensuite, prendre un avis juridique indépendant avant de saisir le régulateur.

D'abord, que sait-on réellement~? Comparer les cours de valorisation à des sources indépendantes, aux cotations des courtiers ou aux contrôles indépendants de prix de la société, et regarder comment l'écart a évolué dans le temps. Une position peu liquide peut avoir une raison légitime à cet écart, donc noter ce qui pourrait l'expliquer.

Ensuite, poser une question neutre si cela est naturel~: «~Je n'arrive pas à retrouver ce cours de valorisation. Qu'est-ce qui m'échappe~?~» Cela règle souvent la question. Si ce n'est pas le cas, apporter ses faits écrits à son responsable. Si son responsable est impliqué, ou si rien ne change, utiliser le canal propre à la société~: la gestion des risques, le product control, la conformité ou une ligne d'alerte. Conserver une trace datée de ses démarches. Il n'est pas nécessaire d'être certain pour le signaler. Établir la vérité relève de leur travail.

Si l'on a utilisé les canaux de la société et que rien ne se passe toujours, prendre un avis juridique indépendant sur le signalement au régulateur. La protection juridique des lanceurs d'alerte varie d'un pays à l'autre. Ne pas sortir de documents de la société, ni s'adresser à quiconque en dehors d'elle, avant d'avoir obtenu cet avis.

\iqwrong{«~Ce n'est pas ma position, donc cela ne me regarde pas~», ou craindre pour sa prime. La première attitude peut rendre complice si c'est avéré. Un mauvais marquage est le moyen par lequel une perte cachée peut grossir pendant des années. L'erreur inverse consiste à accuser publiquement sur un desk, sur la base d'un seul cours.}
\iqfollow{Et si signaler cela nuisait à votre carrière~?}{C'est possible, et il faut le dire, en précisant qu'on en a conscience. Dire aussi que se taire, sur quelque chose qui pourrait être grave, représente un risque plus grand pour soi comme pour la société.}
\end{solution}

\begin{solution}{iq:iv2:behavioural-and-fit:30}
Rester calme, vérifier une fois à voix haute, et si l'on y croit toujours, le dire poliment, en montrant son raisonnement afin que l'interlocuteur puisse désigner l'étape avec laquelle il n'est pas d'accord. S'il a raison, le dire aussitôt, nommer son erreur et poursuivre.

Souvent, le recruteur teste le candidat. Soit il veut voir si le candidat tient une réponse juste sous pression, soit il a mal entendu, soit il pense à une autre version du problème. Le premier geste utile est donc de demander~: «~Avec quelle étape n'êtes-vous pas d'accord~?~» Ensuite, vérifier un cas simple, ou essayer une autre méthode. Dire ce que l'on trouve.

\textbf{Illustration des formulations.} «~Laissez-moi vérifier cela. J'ai supposé [hypothèse], ce qui donne [étape]. Un petit cas donne aussi [résultat]. J'obtiens donc toujours [réponse]. Ai-je mal lu la question, ou voyez-vous une erreur dans une étape~?~» Et si l'on constate qu'il a raison~: «~Vous avez raison. Mon erreur était [étape]. Cela change la réponse en [réponse].~»

Ces deux issues sont bonnes, et la seconde n'est pas un échec. Abandonner aussitôt une réponse juste que l'on sait défendre est un mauvais résultat. Défendre une réponse que l'on vient de découvrir fausse est tout aussi mauvais.

\iqwrong{Céder à la première fronce de sourcils, ce qui laisse penser que son point de vue ne résistera pas à la pression d'un trader. Se braquer ou argumenter sans montrer le moindre calcul est tout aussi mauvais.}
\iqfollow{Et s'il insiste, alors que vous êtes sûr de vous~?}{Dire que l'on voit les choses autrement, et que l'on est prêt à envoyer une courte démonstration après l'entretien. Passer à autre chose, et ne pas laisser cela gâcher la suite de l'entretien.}
\end{solution}
